%%
% 引言或背景
% 引言是论文正文的开端,应包括毕业论文选题的背景、目的和意义;对国内外研究现状和相关领域中已有的研究成果的简要评述;介绍本项研究工作研究设想、研究方法或实验设计、理论依据或实验基础;涉及范围和预期结果等。要求言简意赅,注意不要与摘要雷同或成为摘要的注解。
% modifier: 黄俊杰(huangjj27, 349373001dc@gmail.com)
% update date: 2017-04-15
%%

\chapter{绪论}
\label{cha:introduction}

\section{选题背景与意义}
\label{sec:background}

云原生技术的普及使现代软件系统呈现出部署规模大、服务依赖深、运行状态动态变化快等特征。相比传统单体应用，微服务架构在可扩展性、持续交付和资源利用率方面具有明显优势，但也将问题定位的复杂度从“单组件错误排查”转变为“跨服务链路推理”。一旦系统出现性能抖动、依赖失联、资源争用或配置漂移，运维工程师往往需要在日志、指标、调用链、Kubernetes 事件和变更记录之间反复切换，才能逐步逼近真实根因。

站点可靠性工程（Site Reliability Engineering, SRE）为这一类复杂系统提供了制度化的方法论。SRE 借助告警、服务等级目标（SLO）、故障演练和标准操作流程等机制来保障系统可靠性。然而，传统 SRE 体系仍然高度依赖人类专家的经验积累与手工排障：一方面，故障排查指南（Troubleshooting Guide, TSG）和标准作业流程（Standard Operating Procedure, SOP）通常以静态文档形式存在，更新速度难以跟上业务系统的快速迭代；另一方面，真实故障场景中往往伴随信息过载、告警噪声和异常耦合，工程师很难在有限时间内从大量遥测数据中准确抽取关键证据。

近年来，大语言模型（Large Language Model, LLM）在长上下文理解、自然语言推理和工具调用方面表现出较强潜力，使其成为构建智能运维代理的重要技术基础。LLM 驱动的 Agent 不仅能够理解告警上下文，还可以基于任务目标自主选择诊断动作，执行脚本命令，整合多源观测结果，并输出结构化的根因分析结论。对于 SRE 场景而言，这种能力意味着系统有机会从“辅助决策工具”进一步迈向“可执行、可反思、可演化”的智能运维主体。

但需要指出的是，直接将通用 LLM 投入生产运维场景仍存在明显鸿沟。首先，云环境中存在强烈的环境特异性，不同业务系统的拓扑结构、命名规范、指标分布和故障模式差异极大，通用模型缺乏足够的领域先验。其次，实际运维任务不仅需要文本推理，还需要与具体基础设施进行稳定交互，例如查询 Pod 状态、抓取 Prometheus 指标、分析 Jaeger 链路以及执行受控故障注入。再次，智能体一旦产生幻觉或执行错误动作，可能放大生产事故风险，因此其自主行为需要被限制在可回滚、可追踪、可评估的边界之内。

基于上述背景，本文尝试从“主动演练”和“持续学习”两个方向切入，设计一种适用于云原生环境的自主演化 SRE 框架。其核心思想是：不再让智能体被动等待线上事故，而是通过离线混沌演练构造故障样本、生成诊断轨迹、提炼结构化知识；当线上诊断遇到失配或失败时，再借助反馈闭环识别知识盲区，触发新的演化周期，从而实现对特定业务环境的逐步适配。该思路对于降低运维对资深专家的依赖、缩短根因定位时间、构建面向云基础设施的自愈能力具有明显的理论意义与工程价值。

\section{国内外研究现状和相关工作}
\label{sec:related_work}

\subsection{传统 AIOps 与根因分析方法}

传统 AIOps 的研究重点多集中在异常检测、故障分类和根因定位等单点任务上，常见方法包括基于规则的诊断系统、依赖图传播、统计异常检测以及监督或半监督学习模型。这类方法在稳定场景下具备较好的效率，但通常依赖预先定义的特征工程和标签体系，难以处理拓扑变化频繁、故障组合复杂或上下文信息不完整的生产环境。当系统发生此前未见的新型故障时，传统方法往往只能给出局部提示，而无法形成跨指标、跨服务的完整排障链路。

\subsection{LLM 驱动的运维智能体}

随着大语言模型推理与工具使用能力的增强，研究者开始尝试利用 LLM 直接参与云事件分析与根因定位。已有工作表明，LLM 可以在云事故文本、日志摘要和告警上下文基础上生成较具可读性的根因分析结果 \cite{chen2024automaticrca}。进一步地，围绕自主云系统构建的研究总结了 Agent 在真实运维场景下面临的关键挑战，包括动作空间受限、上下文稀缺、系统安全边界以及多工具协同等问题 \cite{shetty2024buildingautonomousclouds}。这些研究说明，LLM 具备进入运维流程的潜力，但单纯依赖语言理解并不足以支撑高可靠的自动诊断。

在多智能体协作方面，STRATUS \cite{chen2025stratus} 代表了当前较为完整的端到端自主 SRE 实践。该系统将事故处理拆解为检测、诊断、缓解等多个阶段，并通过角色专门化与事务性安全约束控制自治风险。相比单一智能体模式，多智能体协作更适合处理运维任务中的分工问题，例如信息收集、假设验证、证据归纳与策略反思可以分别由不同角色承担。

\subsection{评测基准与演练环境}

LLM Agent 在 SRE 领域的一大难点，是缺少统一、可复现、可交互的评估环境。ITBench 从真实 IT 自动化任务角度构建了智能体评测基准，为跨任务能力比较提供了参考 \cite{jha2025itbench}。AIOpsLab \cite{chen2025aiopslab} 则更贴近云运维场景，提供了微服务部署、故障注入、工作负载生成、Prometheus 指标导出与 Jaeger 链路采集等完整能力，使研究者能够在统一环境下测试诊断 Agent 的行为链条。此类平台为本文的实验原型提供了重要的实现参考。

\subsection{自适应记忆与上下文演化}

除任务执行本身外，近期研究也开始关注如何让黑盒 LLM 在不进行高成本微调的情况下积累经验。Dynamic Cheatsheet 通过自适应记忆机制增强测试时学习能力 \cite{suzgun2025dynamiccheatsheet}；ACE 则进一步将上下文工程与演化思想引入智能体系统，使模型能够根据反馈不断调整自身的上下文表示 \cite{zhang2025ace}。这类工作启发本文从“知识库持续迭代”角度设计运维智能体，而非将 LLM 仅视作一次性的推理器。

\subsection{现有工作的不足}

综合以上研究，可以看到当前 LLM 赋能 SRE 已具备良好开端，但仍存在以下不足。
\begin{enumerate}
    \item \textbf{环境适配能力不足。} 许多方法在通用数据集或固定基准上表现良好，但进入新的微服务系统后，往往缺乏针对该环境的知识沉淀过程。
    \item \textbf{离线学习与在线诊断割裂。} 现有系统通常强调单次诊断成功率，却较少关注失败样例如何反向驱动新一轮演练与知识修补。
    \item \textbf{故障覆盖层次有限。} 多数方法仍以日志分析或指标检测为核心，对网络故障、部署失配、内核级异常等复杂场景支持不足。
    \item \textbf{可追溯性与可解释性有待增强。} 如果缺少统一的诊断轨迹、知识条目与历史记录结构，后续复盘和系统演进将难以开展。
\end{enumerate}

\section{本文研究内容}
\label{sec:research_content}

针对上述问题，本文围绕 OpsEvolver 原型系统展开研究，主要工作包括以下四个方面。
\begin{enumerate}
    \item 设计一套面向云原生场景的多智能体系统架构，将离线演化、在线诊断和反馈闭环纳入统一工作流之中，使不同智能体能够围绕同一诊断目标协同工作。
    \item 实现覆盖网络、容器、Pod、部署配置和磁盘 I/O 等多个层次的故障注入能力，并结合 Prometheus、Jaeger 与日志接口构建统一的可观测性采集链路。
    \item 设计结构化知识条目、规则库与历史库三类持久化载体，将演练过程中得到的症状、证据、建议动作和反馈信息转化为可检索、可复用、可追踪的知识资产。
    \item 围绕原型系统开展实验分析，从实验对象、故障覆盖、诊断链路、知识闭环和典型案例等角度验证该方法在真实微服务环境中的可行性。
\end{enumerate}

\section{本文创新点}
\label{sec:contribution}

相较于现有工作，本文的创新点主要体现在以下几个方面。
\begin{enumerate}
    \item \textbf{提出演练驱动的自演化 SRE 框架。} 本文将混沌演练、RCA 诊断与知识沉淀整合为统一闭环，使系统具备从故障演练中自主学习的能力。
    \item \textbf{构建多角色协同的运维智能体体系。} 通过探索智能体、诊断智能体、反思智能体和反馈分析模块的职责分离，提升了系统在复杂故障场景下的信息组织能力与决策可解释性。
    \item \textbf{实现面向环境特化的知识累积机制。} 不依赖高成本微调，而是通过结构化知识条目、规则库和失败分析结果持续修正上下文，使系统对特定业务环境逐步形成领域适配能力。
    \item \textbf{实现多层次故障注入与可追溯存储。} 将 Chaos Mesh 资源级故障、部署失配以及基于 eBPF 的磁盘 I/O 异常纳入统一故障动作框架，并将每轮演化的注入、告警、轨迹和知识条目进行持久化记录。
\end{enumerate}

\section{论文结构与章节安排}
\label{sec:arrangement}

本文共分为五章，各章内容安排如下。

第一章为绪论，介绍研究背景、相关工作、研究内容与创新点，明确本文的研究问题与技术路线。第二章为系统设计，给出 OpsEvolver 的总体架构、关键模块划分、多智能体协作机制以及知识自演化逻辑。第三章为实现细节，从故障注入、诊断动作、遥测采集、告警生成、知识库持久化与场景部署等角度说明原型系统的工程实现。第四章为实验与分析，围绕实验对象、能力覆盖、机制分析和典型案例展开验证与讨论。第五章为总结与展望，对全文工作进行归纳，并给出当前原型的不足以及后续可能的改进方向。

